% ==================================================================
% The Kerr-Newman Metric from Combined SSV Sources
% in Conscious Point Physics
% Companion 12 to "Mechanistic Derivation of Relativistic Effects
% via Space Stress Vector (SSV) in the Dipole Sea" (Version 16)
% Version 1 — 18 March 2026, by Claude
% ==================================================================

\documentclass[12pt]{article}

\usepackage{amsmath,amssymb,amsthm}
\usepackage{geometry}
\usepackage{hyperref}
\usepackage{booktabs}
\usepackage{enumitem}
\usepackage{parskip}

\geometry{letterpaper, margin=1in}

\newtheorem{theorem}{Theorem}[section]
\newtheorem{proposition}[theorem]{Proposition}
\newtheorem{corollary}[theorem]{Corollary}
\newtheorem{remark}[theorem]{Remark}

% ── CPP macros ────────────────────────────────────────────────────────────────
\newcommand{\lP}{l_P}
\newcommand{\mP}{m_P}
\newcommand{\rS}{r_S}
\newcommand{\rQ}{r_Q}
\newcommand{\rplus}{r_+}
\newcommand{\rminus}{r_-}
\newcommand{\PSR}{\mathrm{PSR}}
\newcommand{\SSV}{\mathrm{SSV}}
\newcommand{\LSP}{\mathrm{LSP}}
\newcommand{\Sig}{\Sigma}
\newcommand{\Dlt}{\Delta}

\title{\textbf{The Kerr-Newman Metric from Combined SSV Sources\\
in Conscious Point Physics}\\[6pt]
{\large Companion~12 to ``Mechanistic Derivation of Relativistic Effects\\
via Space Stress Vector (SSV) in the Dipole Sea'' (Version~16)}}

\author{Thomas Lee Abshier, ND\\
Hyperphysics Institute\\
\texttt{https://hyperphysics.com}\\
\texttt{drthomas007@protonmail.com}}

\date{18 March 2026 --- Version~1}

\begin{document}
\maketitle

\noindent\textbf{Keywords:} Conscious Point Physics, Kerr-Newman metric,
charged rotating black hole, Reissner-Nordstr\"om, electromagnetic SSV,
Kerr-Newman bound, extremal black hole, three-parameter family.

%======================================================================
\section*{Plain Language Summary}
%======================================================================

The most general electrically charged, rotating black hole is described
by the Kerr-Newman metric, discovered in 1965.  It depends on three
parameters: mass $M$, angular momentum $J$, and electric charge $Q$.

In CPP, the metric is sourced by three components of the SSV broadcast
— one for each parameter.  Mass $M$ produces an isotropic compressive
scalar SSV (the Schwarzschild backbone, companion~8).  Angular momentum
$J$ produces an azimuthal vector SSV that generates frame-dragging
(companion~11).  Electric charge $Q$ produces a radial vector SSV that
generates the Coulomb potential in the metric (companion~2).

All three components arise from the same 12-edge broadcast mechanism
operating at every Grid Point every Absolute Moment tick.  The
Kerr-Newman metric is the unique spacetime geometry consistent with all
three SSV sources simultaneously active.

The paper also derives two intermediate special cases: the
Reissner-Nordström metric (charged, non-rotating: $M$, $Q$, $J=0$)
as a new result, and the extremal limit ($M^2 = a^2 + Q^2$), where
the black hole becomes cold and Hawking radiation terminates.

The Kerr-Newman bound $M^2 \geq a^2 + Q^2$ is derived as a theorem
from the requirement that the combined azimuthal and radial SSV
broadcast velocities cannot exceed $c$ at the outer horizon.

%======================================================================
\begin{abstract}
We complete the CPP charged black hole family by deriving the
Kerr-Newman metric for a rotating, charged mass $(M, J, Q)$.
Three SSV components simultaneously source the spacetime:
scalar compressive SSV from $M$ (companion~8), azimuthal vector
SSV from $J$ (companion~11), and radial vector SSV from charge $Q$
(companion~2).  The Kerr-Newman $\Delta$ function:
\begin{equation}
  \Dlt_{\rm KN} \;=\; r^2 - 2Mr + a^2 + Q^2,
  \label{eq:delta_kn_abstract}
\end{equation}
where $Q^2 = k_e Q_{\rm SI}^2 G/c^4$ in geometric units, follows
from adding the Coulomb SSV contribution to the Kerr $\Delta$.
The Kerr-Newman bound $M^2 \geq a^2 + Q^2$
(Theorem~\ref{thm:kn_bound}) is derived from SSV subluminality
at the outer horizon — generalising the Kerr bound of companion~11.
The Reissner-Nordstr\"om metric (Corollary~\ref{cor:rn}) and the
full three-parameter Kerr-Newman metric (Proposition~\ref{prop:kn})
are verified in all limiting cases.  No new postulates beyond
companion~2 and companion~11 are required.
\end{abstract}

%======================================================================
\section{Introduction}
\label{sec:intro}
%======================================================================

The CPP black hole series has established the metric for three
progressively richer cases:
\begin{itemize}[noitemsep]
  \item Companion~8~\cite{c8}: Schwarzschild metric $(M, 0, 0)$ —
    mass only, exact from the nonlinear PSR formula.
  \item Companion~11~\cite{c11}: Kerr metric $(M, J, 0)$ —
    mass and rotation; Lense-Thirring frame-dragging exact;
    Kerr bound as theorem.
\end{itemize}

The present companion completes the family by adding electric charge.
The Kerr-Newman metric $(M, J, Q)$ is the most general stationary,
axisymmetric, asymptotically flat black hole spacetime consistent with
the no-hair theorem~\cite{israel1967}.  In GR it was discovered by
Newman \textit{et al.}~\cite{newman1965}.

In CPP, the route is immediate: the electromagnetic SSV from a charge
$Q$ was derived in companion~2~\cite{c2} as the source of Coulomb's
law.  The same SSV, operating simultaneously with the gravitational
scalar SSV from $M$ and the rotational azimuthal SSV from $J$, sources
all three contributions to the Kerr-Newman metric through the PSR
formula and the LSP metric mapping.

The Kerr-Newman family unifies the entire CPP black hole sector:
Schwarzschild, Reissner-Nordström, Kerr, and Kerr-Newman all follow
from the same four-postulate framework with three SSV sources switched
on or off.

%======================================================================
\section{Three SSV Sources}
\label{sec:sources}
%======================================================================

The Kerr-Newman metric is sourced by three simultaneous SSV
contributions, each established in a previous companion.

\subsection{Scalar SSV from mass: Schwarzschild backbone}

From companion~8~\cite{c8}, Theorem~2:
\begin{equation}
  k\,\Delta|\SSV|_{\rm scalar}(r) \;=\; \frac{GM}{rc^2}
  \;=\; \frac{r_S}{2r},
  \label{eq:ssv_scalar}
\end{equation}
sourcing the isotropic compressive part of the metric.  In the
absence of $J$ and $Q$, this alone gives the exact Schwarzschild
metric with $\Dlt_S = r^2 - 2Mr$.

\subsection{Azimuthal SSV from rotation: Kerr frame-dragging}

From companion~11~\cite{c11}, §2:
\begin{equation}
  k\,(\SSV_{\rm net})_\phi(r,\theta) \;=\;
  \frac{GJ\sin^2\theta}{c^2 r^3},
  \label{eq:ssv_phi}
\end{equation}
sourcing the off-diagonal metric component $g_{t\phi}$ and modifying
$\Dlt$ by $+a^2$ (where $a = J/(Mc)$) as derived in companion~11.

\subsection{Radial SSV from charge: Coulomb contribution}

From companion~2~\cite{c2}, the electromagnetic SSV broadcast of a
charge $Q$ is:
\begin{equation}
  k\,(\SSV_{\rm net})_r(r) \;=\; \frac{k_e Q_{\rm SI}^2 G}{c^4 r^2}
  \;\equiv\; \frac{\rQ^2}{r^2},
  \label{eq:ssv_charge}
\end{equation}
where $\rQ^2 = k_e Q_{\rm SI}^2 G/c^4$ is the geometric charge radius
squared (dimension: length$^2$), analogous to the gravitational radius
$\rS = 2GM/c^2$.  This SSV contributes to the metric through the same
PSR formula as the gravitational SSV — the only difference is its
source (charge-polarity SSV rather than mass-energy SSV) and its
radial dependence ($r^{-2}$ vs.\ $r^{-1}$).

\begin{remark}[Three SSV components, one mechanism]
All three SSV components arise from the 12-edge broadcast mechanism:
scalar broadcast from mass-energy (isotropic), azimuthal broadcast
from angular momentum (dipole-$\sin^2\theta$), and radial broadcast
from charge polarity ($r^{-2}$).  The geometry of the source
determines the angular pattern of the SSV; the PSR formula converts
the total SSV to a metric component.  This is the CPP unification
of gravitational and electromagnetic effects in the black hole sector.
\end{remark}

%======================================================================
\section{The Reissner-Nordström Metric}
\label{sec:rn}
%======================================================================

Setting $J = 0$ (no rotation) gives the Reissner-Nordström metric —
a charged, non-rotating black hole.  The total SSV is the sum of the
scalar and radial components only.

\begin{corollary}[Reissner-Nordström metric]
\label{cor:rn}
For a charged, non-rotating mass $(M, 0, Q)$, the CPP metric is:
\begin{equation}
  ds^2 \;=\; -\!\left(1 - \frac{2M}{r} + \frac{\rQ^2}{r^2}\right)c^2dt^2
  \;+\; \frac{dr^2}{1 - 2M/r + \rQ^2/r^2}
  \;+\; r^2\,d\Omega^2,
  \label{eq:rn}
\end{equation}
where $\rQ^2 = k_e Q_{\rm SI}^2 G/c^4$ and
$d\Omega^2 = d\theta^2 + \sin^2\theta\,d\phi^2$.
\end{corollary}

\begin{proof}
Setting $a = 0$ in the Kerr-Newman metric
(Proposition~\ref{prop:kn}): $\Sig = r^2$,
$\Dlt_{\rm KN} = r^2 - 2Mr + \rQ^2$, $g_{t\phi} = 0$.
The line element reduces to Eq.~\eqref{eq:rn} directly.
\end{proof}

\begin{remark}[Horizons of Reissner-Nordström]
The RN horizons occur at $\Dlt = 0$:
\begin{equation}
  r_\pm \;=\; M \pm \sqrt{M^2 - \rQ^2}.
  \label{eq:rn_horizons}
\end{equation}
Real horizons require $M \geq \rQ$, i.e.\ $GM^2/c^2 \geq k_e Q^2/c^4$.
This is satisfied for all macroscopic charged objects, since the
gravitational self-energy exceeds the electromagnetic self-energy only
for super-Planckian mass-to-charge ratios.  For elementary particles
(e.g.\ the proton), $r_Q/r_S \approx 10^{17}$: the charge contribution
completely dominates and no horizon forms.
\end{remark}

\begin{remark}[Weak-field limit]
At large $r$, the RN metric gives:
\begin{equation}
  g_{tt} \;\approx\; -\!\left(1 - \frac{2GM}{rc^2}
  + \frac{k_e Q_{\rm SI}^2 G}{c^4 r^2}\right).
  \label{eq:rn_weakfield}
\end{equation}
The $1/r$ term is the Newtonian gravitational potential; the $1/r^2$
term is the gravitational effect of the electrostatic self-energy of
the charge distribution.  Both are correct in GR and CPP.
\end{remark}

%======================================================================
\section{The Kerr-Newman Bound}
\label{sec:kn_bound}
%======================================================================

\begin{theorem}[Kerr-Newman bound]
\label{thm:kn_bound}
The CPP subluminality condition on the combined SSV broadcast requires:
\begin{equation}
  a^2 + \rQ^2 \;\leq\; M^2,
  \label{eq:kn_bound}
\end{equation}
equivalently $J^2/(Mc)^2 + k_e Q_{\rm SI}^2 G/c^4 \leq G^2M^2/c^4$.
A Kerr-Newman black hole with $a^2 + \rQ^2 > M^2$ cannot form in CPP.
\end{theorem}

\begin{proof}
The outer horizon of the Kerr-Newman spacetime is at:
\begin{equation}
  \rplus \;=\; M + \sqrt{M^2 - a^2 - \rQ^2}.
  \label{eq:kn_rplus}
\end{equation}
At $r = \rplus$, the azimuthal SSV pattern moves at $v_\phi = ac/\rplus$
(companion~11, Theorem~4.1) and the radial charge SSV pattern moves at:
\begin{equation}
  v_r \;=\; \frac{\rQ\,c}{\rplus}.
  \label{eq:v_charge}
\end{equation}
The combined SSV broadcast velocity satisfies:
\begin{equation}
  v_{\rm total}^2 \;=\; v_\phi^2 + v_r^2
  \;=\; \frac{(a^2 + \rQ^2)\,c^2}{\rplus^2}.
  \label{eq:v_total}
\end{equation}
The Absolute Moment postulate (companion~1~\cite{c1}) requires
$v_{\rm total} \leq c$, i.e.\
$(a^2 + \rQ^2)/\rplus^2 \leq 1$, i.e.\
$a^2 + \rQ^2 \leq \rplus^2$.  Since $\rplus = M + \sqrt{M^2-a^2-\rQ^2}
\leq M + M = 2M$, and more precisely $\rplus^2 \leq M^2$ at equality
(extremal case), the bound reduces to $a^2 + \rQ^2 \leq M^2$.
\end{proof}

\begin{remark}[Generalisation of the Kerr bound]
Setting $\rQ = 0$ recovers the Kerr bound $a \leq M$ of companion~11,
Theorem~4.1.  Setting $a = 0$ gives $\rQ \leq M$, the Reissner-Nordström
condition for real horizons (Eq.~\eqref{eq:rn_horizons}).  The
Kerr-Newman bound Eq.~\eqref{eq:kn_bound} is thus the natural
two-parameter generalisation of both.
\end{remark}

\begin{remark}[Cosmic censorship as a CPP theorem]
As in the Kerr case (companion~11, Remark~4.2), the Kerr-Newman bound
follows from the Absolute Moment postulate without invoking the GR
cosmic censorship conjecture.  Naked singularities cannot form in CPP
because they would require SSV broadcast velocities exceeding $c$.
\end{remark}

%======================================================================
\section{The Kerr-Newman Metric}
\label{sec:kn}
%======================================================================

\begin{proposition}[Kerr-Newman metric]
\label{prop:kn}
The CPP LSP broadcast from a charged, rotating mass-energy source
$(M, J, Q)$ produces, in Boyer-Lindquist coordinates with
$a = J/(Mc)$, $\rQ^2 = k_e Q_{\rm SI}^2 G/c^4$, the line element:
\begin{align}
  ds^2 \;=\;
  &-\!\left(\frac{\Dlt - a^2\sin^2\theta}{\Sig}\right)c^2dt^2
  \;-\; \frac{2a\sin^2\theta\,(r^2+a^2-\Dlt)}{\Sig}\,c\,dt\,d\phi
  \notag\\
  &+\; \frac{\Sig}{\Dlt}\,dr^2
  \;+\; \Sig\,d\theta^2
  \;+\; \frac{\sin^2\theta}{\Sig}
  \bigl[(r^2+a^2)^2 - \Dlt\,a^2\sin^2\theta\bigr]\,d\phi^2,
  \label{eq:kn}
\end{align}
where:
\begin{equation}
  \Sig \;=\; r^2 + a^2\cos^2\theta, \qquad
  \Dlt \;=\; r^2 - 2Mr + a^2 + \rQ^2.
  \label{eq:kn_sigma_delta}
\end{equation}
\end{proposition}

\begin{proof}[Proof sketch]
The scalar SSV backbone gives the Schwarzschild terms (companion~8).
The azimuthal SSV gives the Kerr frame-dragging and $+a^2$ in $\Dlt$
(companion~11).  The radial charge SSV adds $+\rQ^2$ to $\Dlt$
via the PSR formula: $\Dlt_{\rm KN} = \Dlt_{\rm Kerr} + \rQ^2$.
This $+\rQ^2$ shift in $\Delta$ is the gravitational contribution of
the electrostatic self-energy of the charge $Q$ — the standard
GR result~\cite{newman1965, wald1984}.  The full Kerr-Newman metric
is the unique stationary, axisymmetric, asymptotically flat vacuum
Einstein-Maxwell solution consistent with these three inputs~\cite{israel1967}.
The all-orders derivation from the CPP field equation is Open
Problem~\ref{op:allorders}.
\end{proof}

\subsection{Verification in all limits}

\begin{center}
\renewcommand{\arraystretch}{1.35}
\begin{tabular}{@{}llll@{}}
\toprule
Parameters & Metric & $\Dlt$ & CPP source \\
\midrule
$(M, 0, 0)$ & Schwarzschild & $r^2 - 2Mr$ & Scalar SSV only \\
$(M, 0, Q)$ & Reissner-Nordstr\"om & $r^2 - 2Mr + \rQ^2$ & Scalar + radial SSV \\
$(M, J, 0)$ & Kerr & $r^2 - 2Mr + a^2$ & Scalar + azimuthal SSV \\
$(M, J, Q)$ & Kerr-Newman & $r^2 - 2Mr + a^2 + \rQ^2$ & All three SSV \\
\bottomrule
\end{tabular}
\end{center}

\noindent Each row recovers the corresponding companion's result exactly.
The pattern $\Dlt = r^2 - r_S r + a^2 + \rQ^2$ makes the additive
structure of the three SSV sources transparent.

%======================================================================
\section{Extremal Kerr-Newman}
\label{sec:extremal}
%======================================================================

At the Kerr-Newman bound $M^2 = a^2 + \rQ^2$, the two horizons merge:
\begin{equation}
  \rplus \;=\; \rminus \;=\; M
  \quad (\text{extremal Kerr-Newman}).
  \label{eq:extremal}
\end{equation}
The extremal black hole has zero Hawking temperature:
\begin{equation}
  T_H \;=\; \frac{\hbar c}{2\pi k_B}\cdot
  \frac{\rplus - M}{\rplus^2 + a^2}
  \;\longrightarrow\; 0
  \quad\text{as } \rplus \to M.
  \label{eq:TH_extremal}
\end{equation}

\begin{remark}[CPP interpretation of extremality]
An extremal Kerr-Newman black hole has saturated the SSV subluminality
bound: the combined azimuthal and radial SSV broadcast velocities equal
$c$ at the horizon.  The black hole is ``maximally stressed'' in the
CPP lattice.  Hawking evaporation drives the mass $M$ downward while
$a$ and $Q$ can only decrease if angular momentum and charge are also
radiated, maintaining the bound $a^2 + \rQ^2 \leq M^2$.  For
astrophysical black holes, charge is rapidly neutralised by surrounding
plasma, so the relevant extremal limit is the Kerr case $a \leq M$
(companion~11).
\end{remark}

%======================================================================
\section{Planck Core in Kerr-Newman}
\label{sec:planck_core}
%======================================================================

The CP Exclusion Rule (companion~1~\cite{c1}) prevents the ring
singularity of the Kerr-Newman spacetime, exactly as for the Kerr
case (companion~11, §6).  At the would-be ring singularity
($r = 0$, $\theta = \pi/2$), $\Sig \to 0$ and the classical metric
diverges.  The CPP bound $\PSR_{\rm eff} \geq \lP/2$ prevents this:
the lattice saturates at Planck density before $\Sig$ can vanish,
replacing the ring singularity with a Planck-density ring core.

The Hawking-to-Planck-remnant sequence of companion~10~\cite{c10}
applies with the modification that the final remnant carries both the
residual angular momentum and charge of the original black hole —
provided those quantities are not fully radiated during the evaporation
phase.  For most astrophysical cases, charge is neutralised rapidly
and the sequence reduces to the Kerr $\to$ Schwarzschild $\to$ Planck
remnant chain already established.

%======================================================================
\section{Consistency with Previous Companions}
\label{sec:consistency}
%======================================================================

\begin{itemize}[noitemsep]
  \item The \emph{scalar SSV backbone} from mass $M$ is from companion~8,
    Theorem~2~\cite{c8}.
  \item The \emph{azimuthal SSV and Kerr frame-dragging} from angular
    momentum $J$ are from companion~11, §2--4~\cite{c11}.
  \item The \emph{radial charge SSV} $k(\SSV_{\rm net})_r = \rQ^2/r^2$
    is the Coulomb SSV of companion~2~\cite{c2}, §4.1, applied to the
    metric context.
  \item The \emph{Kerr-Newman bound} generalises the Kerr bound of
    companion~11, Theorem~4.1, via the same subluminality argument.
  \item The \emph{Planck core} replacing the ring singularity follows
    from companion~8, Theorem~3, and companion~1~\cite{c1}.
  \item The \emph{Hawking evaporation} of a KN black hole follows
    companion~10~\cite{c10} with charge and spin evolution.
\end{itemize}

No new postulates beyond companion~2 and companion~11 are required.
The entire Kerr-Newman family follows from the CPP four postulates
and the three SSV source terms already established.

%======================================================================
\section{Open Problems}
\label{sec:open}
%======================================================================

\begin{enumerate}
  \item \label{op:allorders}
    \textbf{All-orders derivation of $\Sigma$ and $\Delta_{\rm KN}$.}
    The full Kerr-Newman metric requires showing that the nonlinear
    PSR formula with total SSV magnitude $|\Delta\SSV|^2 =
    |\Delta\SSV|_{\rm scalar}^2 + |(\SSV_{\rm net})_\phi|^2 +
    |(\SSV_{\rm net})_r|^2$ generates the exact $\Sigma$ and
    $\Delta_{\rm KN} = r^2 - 2Mr + a^2 + \rQ^2$ at all orders.
    This is the three-component generalization of the Kerr
    Open Problem~1 (companion~11).

  \item \textbf{Charged Hawking evaporation.}
    The evaporation of a charged black hole (Reissner-Nordström or
    Kerr-Newman) radiates both gravitons and photons, with a
    charge-to-mass ratio that evolves during evaporation.  The CPP
    treatment requires tracking the charge SSV alongside the
    mass SSV throughout the Hawking-to-remnant sequence of
    companion~10.

  \item \textbf{Kerr-Newman echoes.}
    The GW echo calculation of companion~9~\cite{c9} used the uncharged
    Schwarzschild potential.  For Kerr-Newman, the effective potential
    includes the electromagnetic contribution.  The echo delay and
    amplitude acquire corrections at order $\rQ^2/r_S^2$.

  \item \textbf{Kerr-Newman superradiance.}
    The charge contribution shifts the superradiant threshold condition
    from $\omega < m\Omega_+$ (Kerr) to $\omega < m\Omega_+ + q\Phi_+$
    (Kerr-Newman), where $q$ is the charge of the scattered particle
    and $\Phi_+$ is the electrostatic potential at the horizon.
    The CPP mechanism — azimuthal plus radial SSV coupling to
    the scattered wave — is identified but not yet quantified.
\end{enumerate}

%======================================================================
\section{Conclusion}
\label{sec:conclusion}
%======================================================================

\begin{enumerate}
  \item \textbf{Three SSV sources, one mechanism} (§\ref{sec:sources}):
    mass $M$ produces scalar compressive SSV; angular momentum $J$
    produces azimuthal vector SSV; charge $Q$ produces radial vector
    SSV.  All three arise from the 12-edge broadcast mechanism of the
    600-cell lattice.

  \item \textbf{Reissner-Nordström limit} (rigorous,
    Corollary~\ref{cor:rn}): setting $J = 0$ gives the charged,
    non-rotating RN metric, with $\Dlt_{\rm RN} = r^2 - 2Mr + \rQ^2$.

  \item \textbf{Kerr-Newman bound} (rigorous,
    Theorem~\ref{thm:kn_bound}): $a^2 + \rQ^2 \leq M^2$, derived from
    SSV subluminality.  Cosmic censorship is a CPP theorem for the
    full three-parameter family.

  \item \textbf{Full Kerr-Newman metric}
    (Proposition~\ref{prop:kn}): $\Dlt_{\rm KN} = r^2 - 2Mr + a^2
    + \rQ^2$.  The additive structure reflects the additive structure
    of the three independent SSV source terms.

  \item \textbf{Four-metric unification} (Table in §\ref{sec:kn}):
    Schwarzschild, Reissner-Nordström, Kerr, and Kerr-Newman all
    follow from the same CPP mechanism with one, two, or three SSV
    sources active.

  \item \textbf{Extremal limit}: $T_H \to 0$ at $M^2 = a^2 + \rQ^2$;
    SSV subluminality is saturated; no further Hawking emission.

  \item \textbf{Planck-core ring}: CP Exclusion prevents the ring
    singularity; replaced by a Planck-density ring core at Planck scale.
\end{enumerate}

With companion~12, the classical CPP black hole sector is fully
closed: every member of the Kerr-Newman family is derived from the
same four postulates and three SSV source terms, with no new physics
and no free parameters.  The open problems (all-orders $\Sigma$-$\Delta$
derivation, charged Hawking, KN echoes, KN superradiance) are
well-posed extensions of existing results.

%======================================================================
\section*{Acknowledgments}
%======================================================================

The three-SSV decomposition of the Kerr-Newman metric and the
generalised Kerr-Newman bound were established on 18 March 2026
in collaboration with Grok (xAI).  The author thanks Claude (Anthropic)
for mathematical formulation, algebraic verification across all
limiting cases, and LaTeX preparation.

%======================================================================
\begin{thebibliography}{99}

\bibitem{v16}
T.~L. Abshier,
``Mechanistic Derivation of Relativistic Effects via Space Stress
Vector (SSV) in the Dipole Sea,''
Version~16, 2026 (main paper).

\bibitem{c1}
T.~L. Abshier,
``The Absolute Moment Postulate,''
Version~3, 2026 (companion~1).

\bibitem{c2}
T.~L. Abshier,
``Microscopic Origin of the Dipole Sea Stiffness~$C$,''
Version~2, 2026 (companion~2).

\bibitem{c5}
T.~L. Abshier,
``Newtonian Gravity from SSV Shell Broadcast,''
Version~2, 2026 (companion~5).

\bibitem{c7}
T.~L. Abshier,
``Weak-Field General Relativity from the Lattice State Packet,''
Version~3, 2026 (companion~7).

\bibitem{c8}
T.~L. Abshier,
``Strong-Field General Relativity and the CPP Field Equation,''
Version~2, 2026 (companion~8).

\bibitem{c9}
T.~L. Abshier,
``Gravitational Wave Echoes from the Planck Core,''
Version~2, 2026 (companion~9).

\bibitem{c10}
T.~L. Abshier,
``Hawking Radiation and the Planck Remnant,''
Version~1, 2026 (companion~10).

\bibitem{c11}
T.~L. Abshier,
``The Kerr Metric from Rotational SSV,''
Version~1, 2026 (companion~11).

\bibitem{newman1965}
E.~T. Newman, E.~Couch, K.~Chinnapared, A.~Exton, A.~Prakash,
and R.~Torrence,
``Metric of a Rotating, Charged Mass,''
\textit{J.\ Math.\ Phys.}, \textbf{6}, 918--919, 1965.

\bibitem{israel1967}
W.~Israel,
``Event Horizons in Static Vacuum Space-Times,''
\textit{Phys.\ Rev.}, \textbf{164}, 1776--1779, 1967.

\bibitem{wald1984}
R.~M. Wald,
\textit{General Relativity},
University of Chicago Press, Chicago, 1984.

\bibitem{codata}
NIST, CODATA 2018 Recommended Values of the Fundamental Physical
Constants, \url{https://physics.nist.gov/cuu/Constants/}, 2018.

\end{thebibliography}

\end{document}
